\section{沿端口和控制位走一遍代码}

经典主通道使用 \code{0x1F0..0x1F7}，控制端口位于 \code{0x3F6}。
数据寄存器宽 16 位，其余命令块寄存器按 8 位访问。设备与 LBA 高位复用在
drive/head 寄存器，LBA 低 24 位分布在三个端口。写寄存器顺序属于 ATA
协议，不能由通用“结构体映射”替代。

\begin{osgridlongtable}{L{0.16\textwidth}L{0.23\textwidth}L{0.24\textwidth}L{0.25\textwidth}}
\textbf{端口偏移} & \textbf{读语义} & \textbf{写语义} & \textbf{访问宽度}\\ \hline
\endfirsthead
\textbf{端口偏移} & \textbf{读语义} & \textbf{写语义} & \textbf{访问宽度}\\ \hline
\endhead
0 & 数据 & 数据 & 16 位，共 256 次/扇区\\ \hline
1 & 错误 & 特性 & 8 位\\ \hline
2 & 扇区计数 & 扇区计数 & 8 位\\ \hline
3--5 & LBA 低、中、高 & LBA 低、中、高 & 8 位\\ \hline
6 & drive/head & drive/head & 8 位\\ \hline
7 & 状态 & 命令 & 8 位\\ \hline
\end{osgridlongtable}

读取命令前先等待 BSY 清除，选择主从设备与 LBA，写入计数和地址，再写 READ
SECTORS 命令。设备置 BSY 处理请求；准备好数据后清 BSY 并置 DRQ。软件同时
检查 ERR 与 DF，只有 \code{BSY=0, DRQ=1, ERR=0, DF=0} 才能读取数据端口。

\topic{状态轮询要区分忙、就绪、错误和浮空总线}
状态 \code{0xFF} 常表示没有设备响应或总线浮空，不能当作所有能力位都置位。
\code{0x00} 也可能表示设备不存在。探测先选择通道和设备，等待必要延时，
再结合 status、signature 和 IDENTIFY 结果判断。

轮询函数返回强状态而不是布尔值：成功、忙超时、设备错误、设备故障和不存在。
上层日志保留状态寄存器与错误寄存器原值，使 QEMU 与真实控制器差异可诊断。

\topic{LBA28 计算同时受协议宽度与镜像范围约束}
LBA28 最多表达 \(2^{28}\) 个扇区，单扇区 512 字节。计算文件范围时先用
64 位检查：
\[
  \mathrm{byte_offset}=\mathrm{LBA}\times512,\qquad
  \mathrm{byte_end}=\mathrm{byte_offset}+
  \mathrm{sector_count}\times512.
\]
乘法与加法均要在执行前证明不溢出，结果还要小于磁盘镜像或分区长度。硬件
寄存器只接受 28 位，不代表软件内部也应使用狭窄类型。

ATA 的扇区计数字段为 8 位，数值零在部分命令中表示 256 个扇区，不能直接
等同“读取零个”。项目初期每次只读一个扇区，简化协议并让失败位置精确，后续
再增加批量传输。

\topic{Stage 1 镜像需要自描述头而不是裸跳转}
本项目的 Stage 1 描述符已经保存 magic、格式版本、头长度、扇区数、加载段、入口、LBA 和
两级校验固定为磁盘契约。所有字段采用小端精确宽度，解析器逐字段读取，不能把
磁盘字节直接重解释为可能带填充的 C++ 结构体。

\begin{osgridlongtable}{L{0.23\textwidth}L{0.18\textwidth}L{0.43\textwidth}}
\textbf{字段} & \textbf{宽度示意} & \textbf{验证}\\ \hline
\endfirsthead
\textbf{字段} & \textbf{宽度示意} & \textbf{验证}\\ \hline
\endhead
magic & 8 字节 & 精确匹配 \code{OSSTAGE1}\\ \hline
version/header size & 各 16 位 & 当前等于 1 和 28\\ \hline
sector count/flags & 各 16 位 & 1--64 个扇区，标志为零\\ \hline
load segment/entry & 各 16 位 & 形成合法实模式地址和入口\\ \hline
payload LBA & 32 位 & LBA28 且负载不越过 2 MiB 镜像\\ \hline
payload checksum & 16 位 & 覆盖补零后的全部负载字\\ \hline
header checksum & 16 位 & 使 512 字节描述符字和为零\\ \hline
\end{osgridlongtable}

magic 只能识别格式，校验只能发现数据变化，都不能替代范围与权限检查。装载器
先验证完整头和所有区间，再读入最终位置，避免损坏当前栈或页表后才发现错误。

\topic{A20 验证需要可恢复的内存别名实验}
A20 关闭时，地址位 20 被清零，\code{0x000000} 与 \code{0x100000}
发生别名。验证选择两个不会覆盖代码、栈和平台数据的地址，保存原字节，写入
不同模式，再比较是否互相影响，最后恢复。

开启路径可依次尝试快速 A20 端口与键盘控制器协议，每条路径都有超时。直接
写控制端口后继续启动只能证明“发出过命令”，验证读回行为才证明地址线实际
开放。

\topic{当前 Stage 1 还没有在使用特性前执行 CPUID}
当前启动目标固定为 QEMU x86-64。Stage 1 直接设置 CR4.PAE 并访问 EFER，
处理器特性检查要到 Kernel 入口以后才执行。这个顺序能够满足当前机器规格，
却不能安全探测任意实体 CPU：不支持对应 MSR 的处理器会在早期产生 \#GP，
此时 IDT 尚未建立，串口也未必来得及报告原因。

实体机适配需要把检查前移。代码应先用 EFLAGS.ID 确认 CPUID 可用，再查询
最大基本叶和扩展叶，最后检查长模式、PAE 与 MSR。能力存在以后，软件仍可
选择 2 MiB 页或是否启用 NX；检测结果和启用策略是两个步骤。

\topic{GDT 描述符按位定义代码环境}
一个传统段描述符包含 base、limit、type、S、DPL、P、AVL、L、D/B 和 G。
字段跨越多个字节，构造函数用位移与掩码显式编码，并用静态断言验证最终 64 位
数值。长模式代码段要求 L=1、D=0，present 和 executable 必须置位。

\begin{osgridlongtable}{L{0.14\textwidth}L{0.24\textwidth}L{0.48\textwidth}}
\textbf{字段} & \textbf{名称} & \textbf{模式切换中的含义}\\ \hline
\endfirsthead
\textbf{字段} & \textbf{名称} & \textbf{模式切换中的含义}\\ \hline
\endhead
P & present & 描述符可被装载\\ \hline
DPL & privilege & Ring 0 或 Ring 3 访问级别\\ \hline
S/type & 类别与权限 & 代码、数据、可读、可写\\ \hline
L & long & 代码段按 64 位子模式执行\\ \hline
D/B & default/big & 16/32 位默认操作数；L=1 时必须为 0\\ \hline
G & granularity & limit 按字节或 4 KiB 粒度\\ \hline
\end{osgridlongtable}

GDTR 保存表基址与最后有效字节偏移，limit 因此为
\(\mathrm{sizeof(GDT)}-1\)。少减一会扩大有效范围，错误更隐蔽；多减一则
最后描述符无法装载。

\topic{保护模式转换的每条指令都有前置条件}
Stage 1 进入时保持 IF=0，随后装载 GDT，把 CR0.PE 置位，再执行远跳转
重装 CS。写 CR0 以后 PE 已经为 1；远跳转再从 GDT 取得 base、limit 和
D=1 属性，使目标位置开始采用 32 位默认解码。远跳转后的代码用
\code{BITS 32} 编码，随后装载数据段和 32 位栈。汇编器 BITS 切换必须与
运行时控制流到达的新模式对齐。

在 PE 置位与远跳转之间不能插入依赖旧分段语义的复杂代码。错误 GDT、无效
选择子或错误编码会产生异常；IDT 尚未可靠时，调试依赖 QEMU 异常日志和
控制寄存器跟踪。

\topic{初始四级页表的每个表项都要可推导}
PML4、PDPT、PD 和 PT 都要求 4 KiB 对齐。表项低位保存 P、RW、US、PWT、
PCD、A、D、PS、G 与软件位，高位保存物理页号和 NX。地址字段必须屏蔽低
12 位，不能把标志混入物理地址。

初始 identity map 让线性地址等于物理地址，使打开分页前后的当前 RIP、栈和
页表仍可访问。若内核最终链接到高半地址，还要同时建立高半映射；跳入高半后
才能移除不需要的低端别名。

使用 2 MiB 大页时，PD 项设置 PS 并直接保存大页物理基址，地址需 2 MiB
对齐。它减少初始表数量，但权限粒度也变粗。代码、只读数据和可写数据需要不同
权限时再拆成 4 KiB 页。

\topic{进入长模式的控制寄存器顺序}
可验证的顺序为：

\begin{enumerate}[label=\arabic*.]
  \item 建立有效 GDT、32 位栈和初始页表。
  \item 在 CR4 设置 PAE，再把 PML4 物理地址写入 CR3。
  \item 通过 RDMSR/WRMSR 在 EFER 设置 LME，按能力选择 NXE。
  \item 在 CR0 设置 PG，同时保持 PE。
  \item 执行远跳转，装载 L=1 的 64 位代码段。
  \item 在 64 位入口重新装载数据段约定、栈和参数。
\end{enumerate}

EFER.LMA 是只读激活结果，不由软件直接写。只有 LME、PAE 与分页组合满足后，
处理器才置 LMA。检查控制寄存器读回值能区分“请求设置”与“模式已激活”。

\topic{ELF64 解析从文件头建立层层边界}
装载器先确认 magic、64 位类别、小端序、当前版本、x86-64 机器类型与可执行
文件类型。程序头表范围计算为：
\[
  [e\_phoff,\ e\_phoff+e\_phnum\times e\_phentsize).
\]
乘法与加法先检查溢出，entry size 至少覆盖项目读取字段，整个范围必须落在
已加载文件内。

对每个 PT\_LOAD，验证文件区间、内存区间、\code{filesz <= memsz}、
对齐关系、权限和目标保留区。所有段先通过验证再复制；重叠段若格式策略不允许
就拒绝，不能依赖遍历顺序决定最终字节。

\topic{BootInfo 是带版本的跨阶段 ABI}
当前 BootInfo 位于物理地址 \code{0x14000}，在恒等映射下也是内核看到的
虚拟地址。版本 2 由十三个连续的 64 位小端字段组成：magic、版本、总长度、原始
Kernel 文件地址与长度、入口、加载段数、CR3 根、恒等映射大小和初始栈顶。
最后三个字段给出物理内存图地址、条目数和 24 字节条目宽度。固定总长度为
104 字节，magic 的内存字节顺序拼成 \code{OSBOOT64}。

这里不使用 \code{size\_t}、裸枚举或宿主指针。磁盘工具可能运行在 ARM64，
目标内核却是 x86-64；只有 \code{uint64\_t} 这样的固定宽度才能让两端对同一
偏移达成稳定解释。C++ 结构体通过 \code{static\_assert} 检查总大小，Stage 1
则用命名偏移逐字段写入。

交接寄存器只传一个 BootInfo 地址：System V AMD64 ABI 规定第一个整数或
指针参数位于 RDI。Stage 1 先把 RSP 切换到 \code{0x03FFF000}，再执行
\code{CALL}；这样内核入口看见的栈状态来自正常函数调用，不依赖跳转前的偶然
对齐。对象、栈、入口和页表均位于低 64 MiB 恒等映射中。

内核收到指针后仍逐字段验证，不能把“调用者是自己的 Stage 1”当作免检理由。
版本或长度不匹配时输出 \code{BOOT\_INFO\_INVALID} 后停止；CR3 读回值还必须
与结构中的页表根相同。未来字段增加时应提升版本或明确旧长度兼容规则，不能
悄悄改变已有偏移。
